\chapter{Wąska całka Denjoya}
\providecommand{\atlasref}[2]{#2}
\label{ma1:ch:waska-calka-denjoya}
\begingroup
\providecommand{\osc}{\operatorname{osc}}
\providecommand{\Var}{\operatorname{Var}}
\providecommand{\mesh}{\operatorname{mesh}}
\providecommand{\supp}{\operatorname{supp}}
\providecommand{\AC}{\mathrm{AC}}
\providecommand{\BV}{\mathrm{BV}}
\providecommand{\HK}{\mathrm{KH}}
\providecommand{\ind}[1]{\chi_{#1}}
\providecommand{\Dn}{\underline D}
\providecommand{\Du}{\overline D}
\providecommand{\KS}{\mathrm{KS}}
Przyjmujemy oznaczenia i konwencje z rozdziału \ref{ma1:ch:kurzweil-henstock}.
\section{Charakterystyka funkcji pierwotnych i wąska całka Denjoya}

Absolutna ciągłość jest zbyt silna, aby obejmować funkcje pierwotne wszystkich pochodnych. Wprowadzamy warunek kontrolujący oscylacje na zbiorach, a następnie jego przeliczalną wersję.
\begin{definicja}[$\AC_*$ i $\AC G_*$]\label{it:def:acgstar}
Niech $F:[a,b]\to\R$ będzie ciągła i $E\subset[a,b]$. Piszemy $F\in\AC_*(E)$, jeżeli dla każdego $\varepsilon>0$ istnieje $\eta>0$ takie, że dla każdej skończonej rodziny przedziałów $[u_j,v_j]$ o parami rozłącznych wnętrzach i końcach $u_j,v_j\in E$ spełniającej $\sum_j(v_j-u_j)<\eta$ mamy
\[
\sum_j\osc_{[u_j,v_j]}F<\varepsilon.
\]
Piszemy $F\in\AC G_*[a,b]$, jeżeli $[a,b]=\bigcup_{n\ge1}E_n$ dla domkniętych zbiorów $E_n$ takich, że $F\in\AC_*(E_n)$. Na całym przedziale $\AC_*([a,b])=\AC[a,b]$: moduł przyrostu nie przekracza oscylacji, a w drugą stronę dla każdego przedziału wybieramy punkty osiągania minimum i maksimum funkcji. Odcinki o tych końcach mają rozłączne wnętrza i sumę długości nie większą niż wyjściowa rodzina. Stosując definicję bezwzględnej ciągłości do tych odcinków, kontrolujemy sumę oscylacji. Natomiast na podzbiorze $E$ warunek oscylacyjny jest silniejszy od kontroli samych przyrostów między punktami $E$.
\end{definicja}
\begin{twierdzenie}[Domknięcie zbioru i działania algebraiczne]\label{it:thm:acstar-closure-algebra}
Dla ciągłej $F:[a,b]\to\R$ i dowolnego $E\subset[a,b]$ zachodzi
$F\in\AC_*(E)$ wtedy i tylko wtedy, gdy $F\in\AC_*(\overline E)$.
W definicji $\AC G_*$ można więc równoważnie używać przeliczalnych pokryć przez dowolne podzbiory.
Klasa $\AC G_*[a,b]$ jest algebrą rzeczywistą z jedynką, zamkniętą także na składanie z zewnętrzną funkcją lipschitzowską na zakresie $F$.
\end{twierdzenie}
\begin{proof}
Kierunek od $\overline E$ do $E$ wynika z inkluzji zbiorów.
Dla kierunku odwrotnego ustalamy $\varepsilon>0$ i wybieramy $\eta>0$ z warunku $\AC_*(E)$ dla błędu $\varepsilon/4$.
Skończoną rodzinę przedziałów o rozłącznych wnętrzach, końcach w $\overline E$ i łącznej długości mniejszej od $\eta/2$ porządkujemy od lewej do prawej, pomijając przedziały jednopunktowe.
Dzielimy ją na dwie grupy według parzystości indeksu. W każdej grupie różne przedziały oddziela dodatnia odległość, ponieważ między nimi leży przedział drugiej grupy o dodatniej długości.
Każdy koniec przybliżamy ciągiem punktów z $E$. Dla dostatecznie dużego wspólnego indeksu przybliżone przedziały mają dodatnie długości, w każdej grupie rozłączne wnętrza, a suma długości jest mniejsza od $\eta$. Suma oscylacji w każdej grupie jest mniejsza od $\varepsilon/4$.
Jeśli końce przedziałów różnią się o co najwyżej $d$, każdy punkt jednego jest oddalony o co najwyżej $d$ od drugiego. Różnica ich oscylacji nie przekracza $2\omega_F(d)$, gdzie $\omega_F(d)=\sup\{|F(x)-F(y)|:x,y\in[a,b],\ |x-y|\le d\}$ jest modułem ciągłości. Przechodząc do granicy w obu skończonych sumach, uzyskujemy łączne oszacowanie $\varepsilon/2<\varepsilon$.
Stąd warunek na $\overline E$; domknięcia zbiorów pokrycia uzasadniają równoważność definicji.

Dla $F,G\in\AC G_*[a,b]$ bierzemy przecięcia $E_n\cap D_m$ zbiorów ich domkniętych pokryć. Jest to przeliczalne pokrycie, na którego każdym zbiorze obie funkcje spełniają $\AC_*$. Dla każdego przedziału $I$ mamy
\[
 \osc_I(F+G)\le\osc_I F+\osc_I G,\qquad
 \osc_I(FG)\le\|F\|_\infty\osc_I G+\|G\|_\infty\osc_I F.
\]
Druga nierówność wynika z tożsamości
$F(x)G(x)-F(y)G(y)=F(x)(G(x)-G(y))+G(y)(F(x)-F(y))$ i supremum po $x,y\in I$.
Normy są skończone dzięki ciągłości. Wybieramy wspólny próg długości dla obu funkcji z błędem $\varepsilon/(2(1+\|F\|_\infty+\|G\|_\infty))$, co kontroluje sumę oscylacji sumy i iloczynu. Mnożenie przez skalar skaluje oscylację przez jego moduł, a funkcje stałe mają oscylację zero.
Dla lipschitzowskiej $\psi$ ze stałą $L$ na zakresie $F$ zachodzi $\osc_I(\psi\circ F)\le L\osc_I F$, więc obowiązuje warunek na tym samym pokryciu. Wszystkie otrzymane funkcje są ciągłe.
\end{proof}
\begin{definicja}[Wąska całka Denjoya]\label{it:def:denjoy-narrow}
Funkcja $f:[a,b]\to\R$ o skończonych wartościach jest całkowalna w wąskim sensie Denjoya, jeśli istnieje ciągła $F\in\AC G_*[a,b]$ taka, że $F'=f$ prawie wszędzie. Wówczas definiujemy $(\mathrm D_*)\int_a^b f=F(b)-F(a)$. Jednoznaczność uzasadnimy we wniosku \ref{it:cor:denjoy-hk}. Na przedziale jednopunktowym przyjmujemy całkę zero. Jeśli dwie funkcje o skończonych wartościach są równe prawie wszędzie, każda pierwotna jednej z nich jest też pierwotną drugiej; po ustaleniu jednoznaczności oznacza to niezmienniczość całki przy zmianach na zbiorze zerowym.
\end{definicja}
\begin{twierdzenie}[Każda różniczkowalna funkcja należy do $\AC G_*$]\label{it:thm:differentiable-acgstar}
Jeśli $F:[a,b]\to\R$ jest ciągła i ma skończoną pochodną w każdym punkcie wnętrza, to $F\in\AC G_*[a,b]$. Nie zakładamy istnienia pochodnych jednostronnych na końcach.
\end{twierdzenie}
\begin{proof}
Dla $n\ge1$ niech $E_n$ będzie zbiorem tych $x$, dla których $|F(y)-F(x)|\le n|y-x|$ dla wszystkich $y\in[a,b]$ z $|y-x|<1/n$. Dla każdego $x\in(a,b)$ ze skończoną pochodną istnieje promień $r>0$, w którym moduły ilorazów różnicowych nie przekraczają $|F'(x)|+1$. Wybierając $n\ge |F'(x)|+1$ i $1/n<r$, dostajemy $x\in E_n$. Zatem $\bigcup_n E_n$ zawiera $(a,b)$, a do pokrycia dołączamy dwa domknięte zbiory jednopunktowe $\{a\},\{b\}$, na których warunek $\AC_*$ jest automatyczny. Każdy $E_n$ jest domknięty: jeżeli $x_k\in E_n$ i $x_k\to x$, to dla każdego ustalonego $y$ z $|y-x|<1/n$ mamy $|y-x_k|<1/n$ dla dostatecznie dużych $k$. W nierówności $|F(y)-F(x_k)|\le n|y-x_k|$ przechodzimy do granicy, korzystając z ciągłości $F$. Tak otrzymujemy warunek definiujący $E_n$ dla każdego dopuszczalnego $y$. Jeśli $u,v\in E_n$ i $v-u<1/n$, to dla każdego $y\in[u,v]$ mamy $|F(y)-F(u)|\le n(v-u)$. Stąd $\osc_{[u,v]}F\le2n(v-u)$. Dla rodziny przedziałów z końcami w $E_n$ i $\sum_j(v_j-u_j)<\min\{1/n,\varepsilon/(2n)\}$ suma oscylacji jest mniejsza od $\varepsilon$.
\end{proof}
\subsection{Kontrola oscylacji na zbiorach zerowych}
Mówimy, że funkcja ciągła $F$ ma zaniedbywalne wahanie
na zbiorze $N\subset[a,b]$, jeżeli dla każdego
$\varepsilon>0$ istnieje funkcja $\delta:N\to(0,\infty)$
taka, że dla każdej skończonej rodziny
$\delta$-drobnych komórek $[u_j,v_j]$ o rozłącznych
wnętrzach, ze znacznikami $t_j\in N\cap[u_j,v_j]$,
zachodzi $\sum_j|F(v_j)-F(u_j)|<\varepsilon$.
Samo pojęcie nie wymaga założenia $\lambda_1(N)=0$
ani istnienia całki, której $F$ jest funkcją pierwotną.

\begin{lemat}[Zaniedbywalne wahanie na zbiorze zerowym]
\label{ma1:kh:lem:acstar-null-variation}\label{it:lem:acstar-null-variation}
Niech $F$ będzie ciągła na $[a,b]$, $E\subset[a,b]$ domknięty i $F\in\AC_*(E)$. Jeżeli $N\subset E$ ma miarę Lebesgue'a zero, to dla każdego $\varepsilon>0$ istnieje cechowanie na $N$, dla którego dowolny skończony drobny podział częściowy $\{(t_j,[u_j,v_j])\}$ ze znacznikami w $N$ spełnia
\[
 \sum_j|F(v_j)-F(u_j)|<\varepsilon.
\]
\end{lemat}
\begin{proof}
Jeśli $E=\varnothing$, również $N=\varnothing$ i suma przyrostów jest zerowa. Dla $a=b$ nie ma komórek o dodatniej długości, więc teza także jest bezpośrednia. Zakładamy dalej $a<b$ i $E\ne\varnothing$. Składowe przedziałowe dopełnienia $E$ w $[a,b]$ są co najwyżej przeliczalne: każda ma dodatnią długość, zawiera liczbę wymierną w swoim wnętrzu, a różne składowe mają rozłączne wnętrza. Niech $J_k$ oznaczają ich domknięcia. Co najwyżej dwie składowe brzegowe nie mają obu końców w $E$. Suma długości wszystkich $J_k$ nie przekracza $b-a$, ponieważ ich wnętrza są rozłączne.

Twierdzimy, że $\sum_k\osc_{J_k}F<\infty$. Wybieramy $\eta_0>0$ z definicji $\AC_*(E)$ dla błędu $1$. Pomijając skończenie wiele składowych, w tym składowe brzegowe, możemy zapewnić, że suma długości pozostałych jest mniejsza niż $\eta_0$. Każda skończona podrodzina pozostałych ma końce w $E$ i spełnia warunek definicji, więc suma jej oscylacji jest mniejsza od $1$. Supremum po podrodzinach daje zbieżność sumy; skończenie wiele pominiętych oscylacji jest skończonych dzięki ciągłości $F$.

Wybieramy skończoną rodzinę $\mathcal J_0$ składowych (w tym brzegowe) tak, by $\sum_{J_k\notin\mathcal J_0}\osc_{J_k}F<\varepsilon/8$. Niech $m$ oznacza jej liczność. Z definicji $\AC_*(E)$ wybieramy $\eta>0$ dla błędu $\varepsilon/4$. Ponieważ $N$ ma miarę zero, z \atlasref{ME-04::thm:lebesgue-outer-regular}{ME-04: regularności zewnętrznej miary Lebesgue'a} istnieje otwarty zbiór $O\supset N$ o mierze mniejszej od $\eta$. Dla $t\in N$ dobieramy $\delta(t)>0$ tak, by jego otoczenie o promieniu $\delta(t)$ leżało w $O$. Dla każdej składowej $J\in\mathcal J_0$, która nie zawiera $t$, dodatkowo wymagamy $\delta(t)<\operatorname{dist}(t,J)$. Jeżeli $t\in J$, jest jego końcem, gdyż $t\in E$; wymagamy wówczas także $\delta(t)<|J|/2$, gdzie $|J|$ oznacza długość tej składowej. Rozpatrujemy tylko składowe dodatniej długości. Z ciągłości $F$ zmniejszamy $\delta(t)$ tak, by $|F(y)-F(t)|<\varepsilon/(8(m+1))$ dla wszystkich $y\in[a,b]$ z $|y-t|<\delta(t)$. Warunki są skończone i wszystkie dopuszczają dodatni promień.

Dla każdej komórki $I_j=[u_j,v_j]$ z rozpatrywanego podziału częściowego zbiór $E\cap I_j$ jest niepusty (zawiera $t_j$) i zwarty. Niech $p_j=\min(E\cap I_j)$, $q_j=\max(E\cap I_j)$. Rozkładamy przyrost:
\[
 |F(v_j)-F(u_j)|
 \le |F(q_j)-F(p_j)|+|F(p_j)-F(u_j)|+|F(v_j)-F(q_j)|.
\]
Przedziały $[p_j,q_j]$ mają rozłączne wnętrza, końce w $E$ i sumę długości nie większą od $\lambda_1(O)<\eta$. Suma pierwszych składników jest więc mniejsza od $\varepsilon/4$. Jeśli któryś z pozostałych odcinków ma dodatnią długość, leży w domknięciu jednej składowej dopełnienia $E$ i przylega do jej końca w $E$. Dla ustalonej składowej mogą wystąpić najwyżej dwa takie odcinki o dodatniej długości, po jednym przy każdym końcu: dwa przy tym samym końcu miałyby wspólne wnętrze. Składowe brzegowe mają co najwyżej jeden taki odcinek.

Dla składowych spoza $\mathcal J_0$ suma tych składników nie przekracza $2\sum_{J_k\notin\mathcal J_0}\osc_{J_k}F<\varepsilon/4$. Jeżeli komórka wchodzi w składową $J\in\mathcal J_0$, to jej znacznik musi należeć do $J$ (w przeciwnym przypadku warunek odległości wyklucza przecięcie). Znacznik jest więc końcem $J$ w $E$. Ponieważ $\delta(t)<|J|/2$, komórka nie może sięgać drugiego końca $J$. Odcinek skrajny w tej luce ma zatem jeden koniec równy znacznikowi i przyrost mniejszy od $\varepsilon/(8(m+1))$. Łącznie składniki z $\mathcal J_0$ wnoszą mniej niż $\varepsilon/4$. Dodanie trzech oszacowań kończy dowód.
\end{proof}

\begin{wniosek}
\label{ma1:kh:cor:acgstar-null-variation}\label{it:cor:acgstar-null-variation}
Jeśli $F\in\AC G_*[a,b]$ i $N\subset[a,b]$ ma miarę zero, to wniosek lematu \ref{ma1:kh:lem:acstar-null-variation} zachodzi dla całego $N$.
\end{wniosek}
\begin{proof}
Niech $[a,b]=\bigcup_{n\ge1}E_n$ będzie domkniętym pokryciem z definicji. Dla $N\cap E_n$ stosujemy lemat z dokładnością $\varepsilon2^{-n-1}$, otrzymując cechowanie $\delta_n$. Każdemu $t\in N$ przypisujemy najmniejsze $n(t)$ z $t\in E_{n(t)}$ i przyjmujemy $\delta(t)=\delta_{n(t)}(t)$. W dowolnym skończonym podziale częściowym grupujemy komórki według $n(t_j)$. Każda grupa spełnia hipotezy odpowiedniego lematu, więc suma modułów przyrostów jest mniejsza od $\sum_{n\ge1}\varepsilon2^{-n-1}=\varepsilon/2$ (dla pustej rodziny suma jest zerowa). W szczególności jest mniejsza od $\varepsilon$.
\end{proof}

\begin{twierdzenie}[Całkowalność pochodnej prawie wszędzie funkcji $\AC G_*$]
\label{ma1:kh:thm:acgstar-derivative-integrable}\label{it:thm:acgstar-derivative-integrable}
Niech $F\in\AC G_*[a,b]$ oraz niech $f:[a,b]\to\R$ ma skończone wartości i spełnia $F'=f$ prawie wszędzie. Wtedy $f\in\HK[a,b]$ i dla każdego $x\in[a,b]$
\[
 F(x)-F(a)=(\HK)\int_a^x f(t)\,\mathrm dt.
\]
\end{twierdzenie}
\begin{proof}
Zakładamy $a<b$; przypadek przedziału jednopunktowego daje całkę zero. Niech $N$ będzie zbiorem miary zero, poza którym istnieje skończona pochodna równa $f$; dołączamy do $N$ oba końce przedziału. Dla $\varepsilon>0$ wybieramy na $N$ cechowanie z wniosku \ref{ma1:kh:cor:acgstar-null-variation} dla błędu $\varepsilon/3$. Funkcja $h=f\chi_N$ jest równa zeru poza $N$ i ma skończone wartości. Lemat \ref{ma1:kh:lem:kh-null-modification} o zmianach na zbiorze miary zero, zastosowany do $|h|$, daje cechowanie dla błędu $\varepsilon/3$. W punktach $N$ bierzemy minimum obu cechowań.

Dla $t\notin N$ z definicji pochodnej dobieramy $\delta(t)>0$, nie większe od cechowania dla $|h|$, aby
\[
 |F(y)-F(t)-f(t)(y-t)|\le
 \frac{\varepsilon}{3(b-a)}|y-t|
 \qquad (|y-t|<\delta(t),\ y\in[a,b]).
\]
Niech $\mathcal P$ będzie dowolnym pełnym podziałem drobnym. Dla komórek ze znacznikami poza $N$ odejmujemy oszacowania dla $y=u_j$ i $y=v_j$; nierówność trójkąta daje błąd nie większy od $\varepsilon(v_j-u_j)/(3(b-a))$. Suma tych błędów nie przekracza $\varepsilon/3$. Dla komórek ze znacznikami w $N$ suma modułów przyrostów $F$ jest mniejsza od $\varepsilon/3$, a suma $|f(t_j)|(v_j-u_j)$ jest mniejsza od $\varepsilon/3$: pełna suma dla $|h|$ jest właśnie tą drugą sumą i ma błąd mniejszy od $\varepsilon/3$ względem całki zero. W konsekwencji
\[
 |S(f,\mathcal P)-F(b)+F(a)|<\varepsilon,
\]
ponieważ suma wszystkich przyrostów $F$ jest teleskopowa. Definicja daje całkowalność i wartość na $[a,b]$. Ten sam argument stosujemy na każdym $[a,x]$: przecięcia zbiorów $E_n$ z podprzedziałem są domkniętym pokryciem, warunek $\AC_*$ pozostaje prawdziwy, a zbiór wyjątkowy pozostaje zerowy po dołączeniu końców. To daje wzór dla wszystkich $x$.
\end{proof}

\subsection{Różniczkowalność prawie wszędzie}
\begin{lemat}[Przedłużenie afiniczne przez luki]
\label{ma1:kh:lem:acstar-affine-extension}\label{it:lem:acstar-affine-extension}
Niech $E\subset[a,b]$ będzie niepustym zbiorem domkniętym, a $F:[a,b]\to\R$ funkcją ciągłą należącą do $\AC_*(E)$. Istnieje funkcja $G\in\AC[a,b]$, równa $F$ na $E$, afiniczna na każdej wewnętrznej składowej dopełnienia $E$ i stała na każdej składowej brzegowej dopełnienia.
\end{lemat}
\begin{proof}
Niech $(l_k,r_k)$ będą wewnętrznymi składowymi dopełnienia $E$; obydwa ich końce należą do $E$. Na $E$ przyjmujemy $G=F$, na $[l_k,r_k]$ kładziemy
\[
 G(x)=F(l_k)+\frac{x-l_k}{r_k-l_k}(F(r_k)-F(l_k)).
\]
Na $[a,\min E]$ i $[\max E,b]$ przyjmujemy wartości stałe, odpowiednio $F(\min E)$ i $F(\max E)$. Definicje są zgodne na końcach. Funkcja $G$ jest ciągła: we wnętrzu luk jest afiniczna lub stała, a dla $x\in E$ rozważamy $y\to x$. Jeżeli $y\in E$, korzystamy z ciągłości $F$. Jeżeli $y$ należy do wewnętrznej luki długości mniejszej od $\eta$, jej oba końce są oddalone od $x$ o mniej niż $|x-y|+\eta$, a $G(y)$ jest kombinacją wypukłą ich wartości. Jednostajna ciągłość $F$ pozwala uczynić różnicę $|G(y)-F(x)|$ dowolnie małą przez wybór $\eta$ i $|x-y|$. Luk długości co najmniej $\eta$ jest skończenie wiele; przy $y\to x$ możemy spośród nich rozpatrywać tylko luki, których domknięcia zawierają $x$, i na nich ciągłość w końcu wynika bezpośrednio ze wzoru afinicznego. To samo dotyczy co najwyżej dwóch luk brzegowych.

Z dowodu lematu \ref{ma1:kh:lem:acstar-null-variation} wiemy, że $\sum_k\osc_{[l_k,r_k]}F<\infty$. W szczególności funkcja
\[
 w(t)=\sum_k\frac{|F(r_k)-F(l_k)|}{r_k-l_k}
             \chi_{(l_k,r_k)}(t)
\]
jest nieujemna, mierzalna i całkowalna w sensie Lebesgue'a: z twierdzenia \atlasref{ME-08::thm:sum-nonnegative-integrals}{ME-08} jej całka jest sumą $|F(r_k)-F(l_k)|$, ograniczoną przez sumę oscylacji. Jeśli $[u,v]$ leży w domknięciu jednej luki, to $|G(v)-G(u)|=\int_u^v w\,\mathrm d\lambda_1$; w luce brzegowej obie strony są zerowe.

Dla dowolnego $[u,v]$ przecinającego $E$ oznaczamy $p=\min(E\cap[u,v])$, $q=\max(E\cap[u,v])$. Przedziały skrajne $[u,p]$ i $[q,v]$ leżą w domknięciach luk lub są jednopunktowe. Jeśli $p=q$, środkowy przyrost wynosi zero, a oba odcinki skrajne mają rozłączne wnętrza. Zatem
\[
 |G(v)-G(u)|\le |F(q)-F(p)|
       +\int_{[u,p]\cup[q,v]} w\,\mathrm d\lambda_1.
\]
Rozważmy teraz skończoną rodzinę przedziałów o rozłącznych wnętrzach. Dla tych przecinających $E$ środkowe przedziały $[p,q]$ także mają rozłączne wnętrza i końce w $E$. Przy dostatecznie małej łącznej długości całej rodziny suma $|F(q)-F(p)|$ jest mniejsza od $\varepsilon/2$ przez $\AC_*(E)$. Wszystkie odcinki skrajne i przedziały leżące w lukach mają rozłączne wnętrza, a ich suma długości nie przekracza sumy długości całej rodziny. Bezwzględna ciągłość całki Lebesgue'a funkcji $w\in L^1$ (lemat \ref{ma1:kh:lem:integral-ac}, wynik ME-17) daje dla dostatecznie małej tej sumy łączny wkład całek mniejszy od $\varepsilon/2$. Otrzymujemy warunek bezwzględnej ciągłości $G$.
\end{proof}

\begin{twierdzenie}[Różniczkowalność prawie wszędzie na zbiorach $\AC_*$]
\label{ma1:kh:thm:acstar-ae-differentiable}\label{it:thm:acstar-ae-differentiable}
Jeśli $F$ jest ciągła na $[a,b]$ i $F\in\AC_*(E)$ dla domkniętego $E\subset[a,b]$, to $F$ ma skończoną pochodną prawie wszędzie na $E\cap(a,b)$. W szczególności każda funkcja $F\in\AC G_*[a,b]$ jest różniczkowalna prawie wszędzie.
\end{twierdzenie}
\begin{proof}
Dla pustego $E$ nie ma czego dowodzić. Niech $G$ będzie przedłużeniem z lematu \ref{ma1:kh:lem:acstar-affine-extension} i $H=F-G$. W przypadku $a=b$ zbiór $E\cap(a,b)$ jest pusty; dalej zakładamy $a<b$. Na $E$ mamy $H=0$. Ponieważ $G$ jest bezwzględnie ciągła, istnieje $g\in\mathcal L^1[a,b]$ z $G(x)=G(a)+\int_a^x g\,\mathrm d\lambda_1$ (twierdzenie \ref{ma1:kh:thm:ac-lebesgue-primitive}). To samo twierdzenie \ref{ma1:kh:thm:ac-lebesgue-primitive} daje $G'=g$ prawie wszędzie. Wystarczy wykazać, że $H'=0$ prawie wszędzie na $E$.

Dla wewnętrznej luki $(l_k,r_k)$ połóżmy $d_k=\osc_{[l_k,r_k]}F$. Z konstrukcji $G$ obie wartości $F(y),G(y)$ leżą pomiędzy minimum i maksimum $F$ na $[l_k,r_k]$. Wobec tego $|H(y)|\le d_k$ dla $y\in[l_k,r_k]$. Ponadto $\sum_k d_k<\infty$ z dowodu lematu \ref{ma1:kh:lem:acstar-null-variation}.

Dla dodatniej liczby całkowitej $m$ rozważmy zbiory
\[
 U_{k,m}=(l_k-md_k,l_k+md_k)\cup(r_k-md_k,r_k+md_k).
\]
Ich miary nie przekraczają $4md_k$. Jeżeli luk jest nieskończenie wiele, zbiór
\[
 B_m=\bigcap_{n\ge1}\bigcup_{k\ge n}U_{k,m}
\]
ma miarę zero, gdyż dla każdego $n$ jego miara nie przekracza $4m\sum_{k\ge n}d_k$, a prawa strona dąży do zera. Przy skończonej liczbie luk przyjmujemy $B_m=\varnothing$. Wykluczamy z $E$ sumę zbiorów $B_m$, wszystkie końce luk (także brzegowych), oraz $a,b$. Jest to zbiór miary zero.

Ustalmy pozostały punkt $x\in E$ i $m\ge1$. Dla wszystkich dostatecznie dużych $k$ punkt $x$ nie należy do $U_{k,m}$. Jeśli $y$ leży w takiej luce, to jeden z jej końców leży pomiędzy $x$ i $y$, ponieważ $x\in E$. Oznaczając go przez $z$, otrzymujemy $|x-y|\ge |x-z|\ge md_k$. Dla $d_k>0$ mamy więc $|H(y)|/|y-x|\le1/m$; dla $d_k=0$ iloraz jest zerowy. Skończenie wiele pozostałych luk, w tym brzegowe, ma domknięcia niezawierające $x$, zatem można wybrać otoczenie $x$, które nie przecina żadnej z nich. Dla $y\in E$ i $y\ne x$ iloraz jest zerowy. W konsekwencji
\[
 \limsup_{\substack{y\to x\\y\ne x}}
 \frac{|H(y)-H(x)|}{|y-x|}\le\frac1m.
\]
Dowolność $m$ daje $H'(x)=0$. Poza zbiorem zerowym, na którym $G'$ nie istnieje, mamy więc $F'=G'$. Dla funkcji $\AC G_*$ stosujemy wynik do każdego domkniętego zbioru $E_n$ z definicji. Przeliczalna suma zbiorów wyjątkowych jest zbiorem miary zero, a zbiory $E_n$ pokrywają przedział.
\end{proof}

\subsection{Kierunek odwrotny: od pierwotnej KH do $\AC G_*$}
\begin{lemat}[Zaniedbywalne wahanie pierwotnej KH]
\label{ma1:kh:lem:kh-null-variation}\label{it:lem:kh-null-variation}
Niech $F(x)=F(a)+(\HK)\int_a^x f$, gdzie $f\in\HK[a,b]$. Dla każdego zbioru zerowego $N$ i $\varepsilon>0$ istnieje cechowanie na $N$, dla którego suma modułów przyrostów $F$ na dowolnym drobnym podziale częściowym ze znacznikami w $N$ jest mniejsza od $\varepsilon$.
\end{lemat}
\begin{proof}
Wybieramy cechowanie dla błędu $\varepsilon/4$ w definicji całki. Lemat Saksa--Henstocka \ref{ma1:kh:lem:kh-saks-henstock} ogranicza sumę modułów różnic $\Delta F_j-f(t_j)(v_j-u_j)$ przez $\varepsilon/2$. Funkcja $|f|\chi_N$ ma skończone wartości i jest zerowa prawie wszędzie, więc jej całka KH wynosi zero (lemat \ref{ma1:kh:lem:kh-null-modification}). Bierzemy minimum wcześniejszego cechowania oraz cechowania dla $|f|\chi_N$ z błędem $\varepsilon/2$. Podział częściowy uzupełniamy przez lemat Cousina \ref{ma1:kh:lem:kh-cousin} na lukach. Z nieujemności funkcji wynika $\sum_j |f(t_j)|(v_j-u_j)<\varepsilon/2$ na pierwotnym podziale częściowym. Dodając oba oszacowania, otrzymujemy tezę.
\end{proof}

\begin{lemat}[Pokrycie z ograniczoną sumą oscylacji]
\label{ma1:kh:lem:kh-bvstar-cover}\label{it:lem:kh-bvstar-cover}
Dla każdej pierwotnej KH $F$ istnieje przeliczalne pokrycie $[a,b]$ zbiorami domkniętymi $E_n$ oraz stałe $C_n<\infty$ takie, że dla każdej skończonej rodziny przedziałów $I_j$ o rozłącznych wnętrzach i końcach w $E_n$ zachodzi $\sum_j\osc_{I_j}F\le C_n$.
\end{lemat}
\begin{proof}
Niech $f$ będzie funkcją podcałkową. Wybieramy cechowanie $\delta$ dla błędu $1$ w definicji całki i definiujemy
\[
 A_n=\{t:\delta(t)>1/n,\ |f(t)|\le n\},\qquad E_n=\overline{A_n}.
\]
Domknięcie bierzemy w $[a,b]$. Z dodatniości cechowania i skończoności $f(t)$ wynika, że zbiory $A_n$ pokrywają przedział. Nie zakładamy ich mierzalności.

Rozważmy najpierw przedziały $[u_j,v_j]$ o długościach dodatnich i mniejszych od $1/(2n)$, z końcami w $E_n$. Wybieramy $z_j$ maksymalizujący $|F(z)-F(u_j)|$ na tym przedziale; $\osc_{[u_j,v_j]}F\le2|F(z_j)-F(u_j)|$. Pomijamy składniki o maksimum zero. Porządkujemy pozostałe przedziały od lewej do prawej i dzielimy na dwie grupy według parzystości indeksu. W każdej grupie różne przedziały oddziela dodatnia odległość, ponieważ pomiędzy nimi leży przedział drugiej grupy o dodatniej długości.

Dla każdego $j$ wybieramy $t_{j,k}\in A_n$ dążące do $u_j$. Ponieważ $z_j>u_j$ oraz $z_j-u_j<1/(2n)$, dla dostatecznie dużego $k$ mamy $0<z_j-t_{j,k}<1/n$. W każdej skończonej grupie, dla wspólnego dostatecznie dużego $k$, przedziały $[t_{j,k},z_j]$ mają rozłączne wnętrza (korzystamy z dodatnich odstępów). Ze znacznikiem $t_{j,k}$ są $\delta$-drobne. Lemat Saksa--Henstocka \ref{ma1:kh:lem:kh-saks-henstock} daje dla każdej grupy
\[
 \sum_j |F(z_j)-F(t_{j,k})|
 \le 2+n\sum_j(z_j-t_{j,k})\le2+n(b-a).
\]
Przechodzimy do granicy dzięki ciągłości $F$. Dodajemy oszacowania obu grup i uwzględniamy czynnik $2$ z nierówności dla oscylacji. Suma oscylacji krótkich przedziałów jest nie większa niż $4(2+n(b-a))$. Przedziałów długości co najmniej $1/(2n)$ jest najwyżej $2n(b-a)$, a każdy ma oscylację nie większą niż $\osc_{[a,b]}F$. Przedziały jednopunktowe nic nie wnoszą. Możemy zatem przyjąć $C_n=4(2+n(b-a))+2n(b-a)\osc_{[a,b]}F$.
\end{proof}

\begin{lemat}[Wzmocnienie kontroli oscylacji do $\AC_*$]
\label{ma1:kh:lem:bvstar-null-to-acstar}\label{it:lem:bvstar-null-to-acstar}
Niech $F$ będzie ciągła, $E\subset[a,b]$ niepusty i domknięty, a sumy oscylacji po skończonych rodzinach przedziałów o rozłącznych wnętrzach i końcach w $E$ niech będą ograniczone wspólną stałą $C$. Jeśli $F$ ma zaniedbywalne wahanie na każdym zbiorze zerowym w sensie lematu \ref{ma1:kh:lem:kh-null-variation}, to $F\in\AC_*(E)$.
\end{lemat}
\begin{proof}
Niech $(l_k,r_k)$ będą wewnętrznymi lukami $E$, a $d_k=\osc_{[l_k,r_k]}F$. Z założenia mamy $\sum_k d_k\le C$. Konstruujemy przedłużenie $G$ funkcji $F|_E$, afiniczne w lukach wewnętrznych i stałe w brzegowych, według wzoru z lematu \ref{ma1:kh:lem:acstar-affine-extension}. Dowód ciągłości tego przedłużenia korzystał tylko z ciągłości $F$ i domkniętości $E$. Funkcja
\[
 w(t)=\sum_k\frac{|F(r_k)-F(l_k)|}{r_k-l_k}\chi_{(l_k,r_k)}(t)
\]
jest nieujemna i całkowalna w sensie Lebesgue'a na mocy twierdzenia \atlasref{ME-08::thm:sum-nonnegative-integrals}{ME-08}, ponieważ $\int_{[a,b]}w\,\mathrm d\lambda_1=\sum_k|F(r_k)-F(l_k)|\le\sum_k d_k\le C$. W każdej luce moduł przyrostu $G$ jest równy całce $w$ po odpowiednim przedziale. Dla przedziału $[u,v]$ przecinającego $E$ oznaczamy $p=\min(E\cap[u,v])$ i $q=\max(E\cap[u,v])$. Mamy
\[
 |G(v)-G(u)|\le |F(q)-F(p)|
 +\int_{[u,p]\cup[q,v]}w\,\mathrm d\lambda_1.
\]
Dla skończonej rodziny przedziałów o rozłącznych wnętrzach środkowe przedziały mają rozłączne wnętrza i końce w $E$, a wszystkie pozostałe odcinki mają rozłączne wnętrza. Zatem suma modułów przyrostów $G$ nie przekracza $C+\int_a^b w\,\mathrm d\lambda_1$. W szczególności $G$ ma ograniczone wahanie.

Wykażemy zaniedbywalne wahanie $G$ na dowolnym zbiorze zerowym $N$. Dla $\varepsilon>0$ z lematu \ref{ma1:kh:lem:integral-ac} wybieramy próg $\eta>0$ dla funkcji $w$ i błędu $\varepsilon/2$. Z regularności zewnętrznej (\atlasref{ME-04::thm:lebesgue-outer-regular}{ME-04}) wybieramy otwarty $O\supset N$ z $\lambda_1(O)<\eta$. Wtedy $\int_{O\cap[a,b]}w\,\mathrm d\lambda_1<\varepsilon/2$. Dla $t\in N$ dobieramy otoczenie cechowania zawarte w $O$. Gdy $t\notin E$, dodatkowo wymagamy, by nie przecinało $E$. Gdy $t\in N\cap E$, bierzemy minimum z cechowaniem zaniedbywalnego wahania $F$ na $N\cap E$ z błędem $\varepsilon/2$. W dowolnym drobnym podziale częściowym ze znacznikami w $N$ rozcinamy komórki przecinające $E$ jak wyżej. W każdym środkowym przedziale $[p,q]$ pozostaje jego znacznik $t\in N\cap E$, gdyż $p\le t\le q$, a przedział jest nadal drobny jako podprzedział pierwotnej komórki; odcinki jednopunktowe pomijamy. Suma modułów środkowych przyrostów jest mniejsza od $\varepsilon/2$ z zaniedbywalnego wahania $F$. Pozostałe odcinki leżą w $O$, mają rozłączne wnętrza i wnoszą nie więcej niż $\int_{O\cap[a,b]}w\,\mathrm d\lambda_1<\varepsilon/2$. Otrzymujemy żądane zaniedbywalne wahanie $G$.

Z ograniczonego wahania $G$ wynika istnienie skończonej pochodnej prawie wszędzie oraz jej całkowalność w sensie Lebesgue'a. Istotnie, z rozkładu Jordana funkcja $G$ jest różnicą dwóch niemalejących funkcji (na przykład z użyciem funkcji wahania $V(x)=\Var_{[a,x]}G$). Ponieważ $G$ jest ciągła, $V$ jest ciągła: wahanie na przedziale malejącym do punktu dąży do zera, co wynika z addytywności wahania, skończonego wahania całkowitego i ciągłości $G$ w tym punkcie. Dokładniej, na $[x,b]$ wybieramy podział, którego suma różni się od wahania o mniej niż $\eta$. Niech $x=x_0<x_1<\cdots<x_m=b$ będzie tym podziałem, a $S$ jego sumą wariacyjną. Dla $x<y<x_1$ podział $y,x_1,\ldots,x_m$ ma sumę co najmniej $S-|G(y)-G(x)|$. Z addytywności wahania (lemat \ref{ma1:kh:lem:variation-additivity}) wynika
\[
 \Var_{[x,y]}G=\Var_{[x,b]}G-\Var_{[y,b]}G
 \le\eta+|G(y)-G(x)|.
\]
Przechodząc z $y\downarrow x$, a następnie z $\eta\downarrow0$, otrzymujemy prawostronną ciągłość $V$. Z lewej wybieramy podział $[a,x]$ i zastępujemy ostatni węzeł $x$ przez $y<x$ leżące za przedostatnim węzłem; tak samo dostajemy $\Var_{[y,x]}G\le\eta+|G(y)-G(x)|$. Na końcach potrzebna jest tylko odpowiednia strona. Przedłużamy obie niemalejące funkcje przez stałe poza $[a,b]$. Wyznaczają skończone dodatnie miary Lebesgue'a--Stieltjesa. Z twierdzenia \atlasref{ME-25::thm:cdf-derivative-25}{ME-25: pochodna funkcji rozkładu} o różniczkowaniu funkcji rozkładu ich pochodne istnieją prawie wszędzie i są gęstościami części bezwzględnie ciągłych tych miar, zatem należą do $L^1$. Niech $g$ będzie różnicą mierzalnych gęstości tych dwóch miar, wybranych ze skończonymi wartościami na całym przedziale. Wówczas $g\in\mathcal L^1[a,b]$ i $g=G'$ prawie wszędzie; do zbioru wyjątkowego dołączamy końce przedziału.

Argument cechowaniowy z dowodu twierdzenia \ref{ma1:kh:thm:acgstar-derivative-integrable} (a nie samo twierdzenie) wymaga wyłącznie zaniedbywalnego wahania na zbiorze wyjątkowym i skończonej pochodnej poza nim; obie własności właśnie udowodniliśmy dla $G$. Stosujemy ten argument do $G,g$ na każdym $[a,x]$ (zaniedbywalne wahanie jest dziedziczone przez podprzedziały). Otrzymujemy $G(x)-G(a)=(\HK)\int_a^x g$. Ponieważ $g\in\mathcal L^1$, ta całka jest równa całce Lebesgue'a: wystarcza kierunek Lebesgue $\Rightarrow$ McShane z twierdzenia \ref{it:thm:mcshane-lebesgue} i inkluzja McShane $\subset\HK$, niezależne od charakterystyki $\AC G_*$. Z \atlasref{ME-25::thm:FTC-lebesgue-25}{ME-25: podstawowego twierdzenia rachunku całkowego} zastosowanego do rzeczywistej funkcji $g\in\mathcal L^1([a,b],\lambda_1)$ wynika $G\in\AC[a,b]$ (dodanie stałej $G(a)$ nie zmienia warunku bezwzględnej ciągłości).

Pozostaje kontrola oscylacji $F$. Gdy $[u,v]$ ma końce w $E$, każda luka wewnętrzna przecinająca $(u,v)$ jest w nim zawarta wraz z końcami w $[u,v]$. Na $E$ mamy $F=G$, a na luce $|F-G|\le d_k$, zatem
\[
 \osc_{[u,v]}F\le\osc_{[u,v]}G+
        2\sum_{[l_k,r_k]\subset[u,v]}d_k.
\]
Wybieramy skończoną rodzinę luk z ogonem $\sum d_k<\varepsilon/4$. Wybieramy $\eta>0$ mniejsze od długości każdej wybranej luki (warunek pomijamy dla pustej rodziny). Jeżeli suma długości rodziny przedziałów z końcami w $E$ jest mniejsza od $\eta$, żaden przedział nie zawiera wybranej luki. Pozostała luka o dodatniej długości może być zawarta w co najwyżej jednym przedziale rodziny, więc łączny drugi składnik jest mniejszy od $\varepsilon/2$. Zmniejszamy $\eta$ tak, by bezwzględna ciągłość $G$ dawała sumę oscylacji mniejszą od $\varepsilon/2$: w każdym przedziale wybieramy punkty minimum i maksimum, a odcinki między nimi mają rozłączne wnętrza i nie większą łączną długość. Dodanie oszacowań daje $\AC_*(E)$.
\end{proof}

\begin{twierdzenie}[Charakteryzacja funkcji pierwotnych KH]
\label{ma1:kh:thm:kh-acgstar-characterization}\label{it:thm:kh-acgstar-characterization}
Dla ciągłej $F:[a,b]\to\R$ równoważne są: (i) $F\in\AC G_*[a,b]$; (ii) istnieje $h\in\HK[a,b]$ takie, że $F(x)=F(a)+\int_a^x h$ dla każdego $x$. Wówczas $F'=h$ prawie wszędzie.
\end{twierdzenie}
\begin{proof}
Z (i) twierdzenie \ref{ma1:kh:thm:acstar-ae-differentiable} daje istnienie skończonej pochodnej prawie wszędzie. Przyjmujemy $h=F'$ w punktach istnienia skończonej pochodnej i zero poza nimi, w tym na końcach. Twierdzenie \ref{ma1:kh:thm:acgstar-derivative-integrable} daje (ii).

Z (ii) lemat \ref{ma1:kh:lem:kh-null-variation} daje zaniedbywalne wahanie na każdym zbiorze zerowym. Lemat \ref{ma1:kh:lem:kh-bvstar-cover} daje przeliczalne domknięte pokrycie $E_n$. Dla każdego niepustego $E_n$ lemat \ref{ma1:kh:lem:bvstar-null-to-acstar} daje $F\in\AC_*(E_n)$, a dla pustego warunek zachodzi automatycznie. Stąd (i). Równość $F'=h$ prawie wszędzie wynika z twierdzenia o różniczkowaniu pierwotnej KH \ref{ma1:kh:thm:kh-primitive-derivative}.
\end{proof}

\begin{wniosek}[Denjoy = KH]\label{it:cor:denjoy-hk}
Wąska całka Denjoya i całka KH mają te same funkcje całkowalne oraz te same wartości. Jeśli $F\in\AC G_*[a,b]$ i $F'=0$ prawie wszędzie, to $F$ jest stała.
\end{wniosek}
\begin{proof}
Dla funkcji KH-całkowalnej jej pierwotna należy do $\AC G_*$ na mocy twierdzenia \ref{ma1:kh:thm:kh-acgstar-characterization} i ma pochodną równą funkcji podcałkowej prawie wszędzie. Odwrotnie, jeśli $F\in\AC G_*$ oraz $F'=f$ prawie wszędzie, to $F-F(a)=\int_a^x h$ dla pewnego $h\in\HK$, z $h=f$ prawie wszędzie. Lemat \ref{ma1:kh:lem:kh-null-modification} o zmianie na zbiorze zerowym pozwala zastąpić $h$ przez $f$. Dla $f=0$ pierwotna jest stała. Jeśli $F_1,F_2$ są dwiema pierwotnymi wąskiej całki Denjoya tej samej funkcji, to z twierdzenia \ref{it:thm:acstar-closure-algebra} wynika $F_1-F_2\in\AC G_*$. Poza sumą zbiorów wyjątkowych mamy $(F_1-F_2)'=0$. Zatem różnica jest stała i przyrosty obu pierwotnych są identyczne, co dowodzi jednoznaczności.
\end{proof}
% Synchronizacja z CAL-08: kryterium na zbiorze wyjątkowym.
\begin{twierdzenie}[Kryterium zaniedbywalnego wahania na zbiorze wyjątkowym]
\label{it:thm:denjoy-exceptional-variation}
Niech $F:[a,b]\to\R$ będzie ciągła, $N\subset[a,b]$
zbiorem miary Lebesgue'a zero zawierającym $a,b$,
a $F$ niech ma skończoną pochodną w każdym punkcie
$[a,b]\setminus N$. Niech $f=F'$ poza $N$, z dowolnymi
skończonymi wartościami na $N$. Następujące warunki
są równoważne:
\begin{enumerate}[label=(\roman*)]
\item $F\in\AC G_*[a,b]$;
\item $F$ ma zaniedbywalne wahanie na $N$ w sensie lematu
\ref{ma1:kh:lem:kh-null-variation};
\item $f\in\HK[a,b]$ i dla każdego $x\in[a,b]$ zachodzi
$F(x)-F(a)=(\HK)\int_a^x f$.
\end{enumerate}
W szczególności warunek (ii) nie wymaga oddzielnego
sprawdzania zaniedbywalnego wahania na wszystkich
zbiorach miary zero.
\end{twierdzenie}
\begin{proof}
Implikacja (i)$\Rightarrow$(ii) wynika z wniosku
\ref{ma1:kh:cor:acgstar-null-variation}.
Dla (ii)$\Rightarrow$(iii) ustalmy $a<x\le b$ i
$\varepsilon>0$. Na $N\cap[a,x]$ ograniczamy cechowanie z warunku (ii)
dla błędu $\varepsilon/3$: każdy drobny podział częściowy
$[a,x]$ jest również podziałem częściowym $[a,b]$.
Stąd to samo oszacowanie zaniedbywalnego wahania
obowiązuje na podprzedziale.
Na całym $[a,x]$ wybieramy także cechowanie $\gamma$
dla funkcji $|f|\chi_N$, której całka KH wynosi zero
z lematu \ref{ma1:kh:lem:kh-null-modification}, z błędem
$\varepsilon/3$. W punktach $N\cap[a,x]$ bierzemy
minimum obu cechowań. W pozostałych punktach $t$
zmniejszamy $\gamma(t)$ tak, aby z różniczkowalności
$F$ wynikało
\[
 |F(y)-F(t)-f(t)(y-t)|
 \le \frac{\varepsilon}{3(x-a)}|y-t|
\]
dla każdego $y\in[a,x]$ w otoczeniu cechowania.
Na dowolnym drobnym podziale $[a,x]$ suma błędów
$|F(v)-F(u)-f(t)(v-u)|$ z komórek ze znacznikiem
poza $N$ nie przekracza $\varepsilon/3$, gdyż
$t\in[u,v]$ i $\sum(v-u)=x-a$.
Dla pozostałych komórek suma modułów przyrostów
$F$ jest mniejsza od $\varepsilon/3$ na mocy (ii),
a suma $|f(t)|(v-u)$ jest mniejsza od $\varepsilon/3$
jako pełna suma dla $|f|\chi_N$.
Teleskopowanie przyrostów daje
$|S(f,\mathcal P)-F(x)+F(a)|<\varepsilon$.
Dla $x=a$ obowiązuje konwencja całki zerowej.
Otrzymujemy (iii). Wreszcie (iii)$\Rightarrow$(i)
jest implikacją z twierdzenia
\ref{ma1:kh:thm:kh-acgstar-characterization}.
\end{proof}

\begin{wniosek}[Zaniedbywalne wahanie na zbiorze przeliczalnym]
\label{it:cor:denjoy-countable-variation}
Każda funkcja ciągła $F:[a,b]\to\R$ ma zaniedbywalne
wahanie na każdym co najwyżej przeliczalnym zbiorze
$N\subset[a,b]$.
\end{wniosek}
\begin{proof}
Dla $N=\varnothing$ teza jest natychmiastowa. W przeciwnym
razie numerujemy jego różne punkty jako $t_k$ (przy zbiorze
skończonym bierzemy tylko skończenie wiele indeksów).
Ustalmy $\varepsilon>0$. Z ciągłości $F$ w $t_k$ wybieramy
$\delta(t_k)>0$ tak, aby $|F(y)-F(t_k)|<
\varepsilon 2^{-k-3}$ dla $y\in[a,b]$ spełniających
$|y-t_k|<\delta(t_k)$.
W drobnym podziale częściowym każda komórka ze znacznikiem
$t_k$ ma przyrost o module mniejszym od
$\varepsilon 2^{-k-2}$ na mocy nierówności trójkąta.
Po pominięciu komórek jednopunktowych co najwyżej dwie
komórki mogą mieć ten sam znacznik: każda zawiera $t_k$,
a ich wnętrza są rozłączne. Gdy występują dwie, punkt
$t_k$ musi być prawym końcem jednej i lewym końcem drugiej.
Stąd suma modułów wszystkich przyrostów jest nie większa
niż $\varepsilon\sum_{k\ge1}2^{-k-1}=\varepsilon/2$.
To jest żądany warunek.
\end{proof}

\begin{przyklad}[Całkowalność z definicji Denjoya]
Dla $F(x)=x^2\cos(\pi/x^2)$, $x>0$, $F(0)=0$, bierzemy domknięte pokrycie $\{0\},[1/2,1],[1/3,1],\dots$ przedziału $[0,1]$. Na każdym $[1/n,1]$ funkcja $F$ jest klasy $C^1$, więc bezwzględnie ciągła, a na zbiorze jednopunktowym warunek $\AC_*$ jest automatyczny. Zatem $F\in\AC G_*$ i $F'$ jest całkowalna w sensie Denjoya z wartością $F(1)-F(0)=-1$. Pochodna nie jest jednak bezwzględnie całkowalna: jest to funkcja z przykładu \ref{ma1:kh:ex:conditional-derivative}, gdzie wykazano rozbieżność całki modułu pochodnej.
\end{przyklad}
\begin{uwaga}[Wąska i szeroka całka Denjoya]
Wąska całka Denjoya ma tę samą klasę funkcji co KH. Szeroka całka Denjoya wykorzystuje słabszą klasę $ACG$ (kontrola przyrostów zamiast oscylacji na zbiorach pokrycia) i pochodną aproksymatywną. Nie utożsamiamy jej z całką KH; dokładną definicję i porównanie podano w części poświęconej szerokiej całce Denjoya poniżej.
\end{uwaga}

% ============================================================
% UZUPELNIENIE MATEMATYCZNE: ROZDZIAL 8
% ============================================================
\subsection{Jak w praktyce rozpoznawać pierwotne wąskiej całki Denjoya}
\begin{twierdzenie}[Przeliczalny zbiór punktów nieróżniczkowalności]\label{ma1:kh:tw:denjoy-countable}\label{tw:denjoy-countable}
Niech $F:[a,b]\to\R$ będzie ciągła oraz różniczkowalna ze skończoną pochodną w każdym punkcie $(a,b)\setminus N$, gdzie $N\subset[a,b]$ jest zbiorem co najwyżej przeliczalnym. Wtedy $F\in\AC G_*[a,b]$. Jeśli funkcja $f:[a,b]\to\R$ ma skończone wartości i $f=F'$ na $(a,b)\setminus N$, to $f$ jest całkowalna w sensie KH oraz w wąskim sensie Denjoya, a obie całki są równe $F(b)-F(a)$.
\end{twierdzenie}
\begin{proof}
Dla $n\ge1$ definiujemy
\[
 E_n=\left\{x\in[a,b]: |F(y)-F(x)|\le n|y-x|
 \ \text{dla wszystkich}\ y\in[a,b]\ \text{z}\ |y-x|<1/n\right\}.
\]
Każdy $E_n$ jest domknięty na mocy argumentu z dowodu twierdzenia \ref{it:thm:differentiable-acgstar}: ścisła nierówność $|y-x|<1/n$ zapewnia, że po przybliżeniu $x$ przez punkty $E_n$ warunek odległości nadal obowiązuje. Każdy punkt, w którym istnieje skończona pochodna (jednostronna na brzegu), należy do któregoś $E_n$. Dla przedziału $[u,v]$ o końcach w $E_n$ i długości mniejszej niż $1/n$ mamy
$\osc_{[u,v]}F\le2n(v-u)$, jak w dowodzie twierdzenia \ref{it:thm:differentiable-acgstar}. Stąd $F\in\AC_*(E_n)$. Zbiory jednopunktowe $\{x\}$, $x\in N\cup\{a,b\}$, są domknięte, a warunek $\AC_*$ zachodzi na nich trywialnie. Przeliczalna rodzina $E_n$ i tych zbiorów jednopunktowych pokrywa $[a,b]$, więc $F\in\AC G_*$. Ponieważ $N\cup\{a,b\}$ ma miarę Lebesgue'a zero, $F'=f$ prawie wszędzie; stosujemy twierdzenie \ref{ma1:kh:thm:acgstar-derivative-integrable} oraz wniosek \ref{it:cor:denjoy-hk}. Alternatywnie cała teza wynika z wniosku \ref{it:cor:denjoy-countable-variation} i twierdzenia \ref{it:thm:denjoy-exceptional-variation}, zastosowanych do $N\cup\{a,b\}$.
\end{proof}
\begin{przyklad}[Funkcja z przeliczalnie wieloma impulsami i punktem skupienia]
Niech $c_n=2^{-n}$, $r_n=2^{-n-3}$ oraz niech $\psi\in C_c^1((-1,1))$ spełnia $\psi(0)=1$. Definiujemy
$F(0)=0$ oraz $F(x)=\sum_{n\ge1}n^{-1}\psi((x-c_n)/r_n)$ dla $x\in(0,1]$. Przedłużamy $\psi$ przez zero poza $(-1,1)$; przedłużenie jest klasy $C^1$, gdyż nośnik jest zwartym podzbiorem $(-1,1)$. Nośnik $n$-tego składnika leży w $[c_n-r_n,c_n+r_n]$. Ponieważ $c_{n+1}+r_{n+1}=9c_n/16<7c_n/8=c_n-r_n$, przedziały te są parami rozłączne i skupiają się tylko w $0$. Szereg jest lokalnie skończony poza $0$, więc $F$ jest tam klasy $C^1$. W otoczeniu $0$ mogą występować wyłącznie składniki o dowolnie dużych indeksach; ich moduły są ograniczone przez $\|\psi\|_\infty/n$, co dowodzi ciągłości w $0$. Niech $f(x)=F'(x)$ dla $x\in(0,1)$ oraz $f(0)=f(1)=0$.
Z twierdzenia \ref{ma1:kh:tw:denjoy-countable}, przy $N=\{0,1\}$,
wynika $f\in\HK[0,1]$ oraz $(\HK)\int_0^1f=F(1)-F(0)=0$.
Dla każdego $n$ wartości $F$ w punktach $c_n-r_n,c_n,c_n+r_n$
wynoszą kolejno $0,1/n,0$. Dla dowolnego $m$ włączamy te trzy
punkty dla $1\le n\le m$ do jednego podziału $[0,1]$.
Ponieważ przedziały nośników są rozłączne, suma modułów
przyrostów na tym podziale wynosi co najmniej
$2\sum_{n=1}^m1/n$. Rozbieżność szeregu harmonicznego
daje $V_0^1(F)=+\infty$. W istocie na $n$-tym przedziale pochodna wynosi $(nr_n)^{-1}\psi'((x-c_n)/r_n)$. Podstawienie $t=(x-c_n)/r_n$ na każdym przedziale
$[c_n-r_n,c_n+r_n]$ jest dopuszczalne w całce Riemanna,
ponieważ $\psi'$ jest ciągła. Zgodność z całką Lebesgue'a
wynika z \atlasref{ME-08::thm:riemann-to-lebesgue}{ME-08}.
Poza sumą tych przedziałów funkcja $f$ jest zerowa.
Z rozłączności przedziałów i twierdzenia
\atlasref{ME-08::thm:sum-nonnegative-integrals}{ME-08} otrzymujemy
\[
 \int_{[0,1]}|f(x)|\,\mathrm d\lambda_1(x)
 =\left(\int_{-1}^1|\psi'(t)|\,\mathrm dt\right)\sum_{n\ge1}\frac1n=\infty.
\]
Stały czynnik jest co najmniej $2$, gdyż $\psi(-1)=\psi(1)=0$, $\psi(0)=1$, a wzór Newtona--Leibniza i nierówność trójkąta na $[-1,0]$ oraz $[0,1]$ dają po co najmniej $1$. Stąd $F\notin\BV[0,1]$ i $F\notin\AC[0,1]$. Mimo to całka KH pochodnej istnieje.
\end{przyklad}
\begin{stwierdzenie}[Sklejanie pierwotnych]\label{ma1:kh:st:denjoy-sklejanie}\label{st:denjoy-sklejanie}
Niech $a<c<b$ i niech $F:[a,b]\to\R$ będzie ciągła. Jeśli $F|_{[a,c]}\in\AC G_*[a,c]$ oraz $F|_{[c,b]}\in\AC G_*[c,b]$, to $F\in\AC G_*[a,b]$. Jeśli dodatkowo $f:[a,b]\to\R$ ma skończone wartości i $F'=f$ prawie wszędzie na obu podprzedziałach, to $(\mathrm D_*)\int_a^b f=(\mathrm D_*)\int_a^c f+(\mathrm D_*)\int_c^b f$.
\end{stwierdzenie}
\begin{proof}
Domknięte pokrycia obu przedziałów przez zbiory, na których obowiązuje $\AC_*$, są także pokryciami przez zbiory domknięte w $[a,b]$. Łącząc te dwie rodziny, otrzymujemy przeliczalne pokrycie $[a,b]$ bez konieczności kontrolowania przedziałów mających końce po różnych stronach $c$: definicja $\AC G_*$ wymaga kontroli osobno na każdym zbiorze pokrycia. Jeżeli $f$ ma skończone wartości i $F'=f$ prawie wszędzie na obu częściach, to równość zachodzi także prawie wszędzie na całym przedziale. Z definicji \ref{it:def:denjoy-narrow} wąskiej całki Denjoya istnieją wszystkie trzy całki, a wzór wynika z $F(b)-F(a)=(F(c)-F(a))+(F(b)-F(c))$.
\end{proof}
\begin{uwaga}[Pochodna prawie wszędzie nie identyfikuje pierwotnej]
Ciągłość i równość $F'=f$ prawie wszędzie nie wystarczają do wzoru Newtona--Leibniza. Przykładem jest funkcja Cantora $C$ skonstruowana w \atlasref{ME-17::ex:cantor-measure-17}{ME-17: przykładzie miary Cantora}, dla której $C'=0$ prawie wszędzie, lecz $C(1)-C(0)=1$. Warunek $\AC G_*$ usuwa tę niejednoznaczność. Jest zatem warunkiem na \emph{pierwotną}, a nie tylko na zbiór punktów istnienia jej pochodnej.
\end{uwaga}

% Import IT-05: zachowane etykiety, odsyłacze do lokalnych wyników AM1.
\begin{twierdzenie}[Ograniczone wahanie pierwotnej Denjoya]\label{it:thm:denjoy-bv-ac}
Zachodzi $\AC G_*[a,b]\cap\BV[a,b]=\AC[a,b]$.
Jeżeli $f$ jest całkowalna w wąskim sensie Denjoya i
$F(x)=(\mathrm D_*)\int_a^x f$, to $F$ ma ograniczone wahanie wtedy i tylko wtedy, gdy $f\in\mathcal L^1[a,b]$.
W takim przypadku $V_a^b(F)=\int_a^b|f|\,\mathrm d\lambda_1$, gdzie $V_a^b$ oznacza wahanie całkowite.
\end{twierdzenie}
\begin{proof}
Absolutna ciągłość pociąga $\AC G_*$ przez pokrycie całym przedziałem, a ograniczone wahanie wynika z twierdzenia \ref{ma1:kh:thm:ac-lebesgue-primitive}.
Odwrotnie, niech $F\in\AC G_*\cap\BV$. Rozkład Jordana przedstawia $F$ jako różnicę dwóch niemalejących funkcji. Ich pochodne istnieją prawie wszędzie i są całkowalne Lebesgue'a, jak w dowodzie lematu \ref{ma1:kh:lem:bvstar-null-to-acstar}. Zatem $F'\in\mathcal L^1$, po nadaniu skończonych wartości na zbiorze wyjątkowym.
Z twierdzenia \ref{ma1:kh:thm:acgstar-derivative-integrable}, zgodności całek McShane'a i Lebesgue'a (twierdzenie \ref{it:thm:mcshane-lebesgue}) oraz inkluzji McShane'a w KH otrzymujemy
\[
 F(x)=F(a)+\int_a^x F'(t)\,\mathrm d\lambda_1(t).
\]
Twierdzenie \ref{ma1:kh:thm:ac-lebesgue-primitive} daje $F\in\AC$.
Dla funkcji $f$ całkowalnej w wąskim sensie Denjoya
wniosek \ref{it:cor:denjoy-hk} daje $f\in\HK[a,b]$
oraz zgodność całek na każdym podprzedziale $[a,x]$,
gdyż ograniczenie pierwotnej klasy $\AC G_*$ pozostaje
w tej klasie. Z twierdzeń
\ref{ma1:kh:thm:kh-acgstar-characterization}
i \ref{ma1:kh:thm:kh-primitive-derivative}
wynika więc $F\in\AC G_*$ i $F'=f$ prawie wszędzie.
Jeżeli $F\in\BV$, to z pierwszej części dowodu
$F\in\AC$, a z twierdzenia
\ref{ma1:kh:thm:ac-lebesgue-primitive} wynika
$f\in\mathcal L^1$. Odwrotnie, dla $f\in\mathcal L^1$
zgodność całek Lebesgue'a i KH daje reprezentację
$F$ jako pierwotnej Lebesgue'a, zatem $F\in\AC\subset\BV$.
Dla $a<b$ twierdzenie \ref{ma1:thm:ac-banach-separable}
z rozdziału 37 daje dokładnie
\[
 V_a^b(F)=\Var_{[a,b]}F
 =\int_{[a,b]}|F'(t)|\,\mathrm d\lambda_1(t)
 =\int_{[a,b]}|f(t)|\,\mathrm d\lambda_1(t),
\]
gdzie ostatnia równość wynika z równości prawie wszędzie.
Dla $a=b$ wahanie i całki są zerowe, a wszystkie
rozważane funkcje należą do odpowiednich klas.
\end{proof}
% Import IT-09; zachowane etykiety źródłowe.
\endgroup
